\documentclass[xelatex,mathserif,serif,notheorems]{beamer} % ,notheorems
\mode<presentation>{
  %% BEAMER CONFIG %%
  \usetheme{Madrid}
  \setbeamertemplate{navigation symbols}{}
  \setbeamersize{text margin left=0.5cm,text margin right=0.5cm}
  \usefonttheme[onlymath]{serif}
}

% \usepackage[latin1]{inputenc}
\usepackage[english]{babel}
% \usepackage{a4}
\usepackage{latexsym}
\usepackage{amssymb}
\usepackage{amsmath}
\usepackage{epsfig}
\usepackage[T1]{fontenc}
% \usepackage{mathptmx} % Use Times New Romans as font?
\usepackage{color}
\usepackage{epstopdf}
\usepackage{microtype}
\usepackage{hyperref}
\usepackage[useregional]{datetime2}
\DTMlangsetup[en-US]{showdayofmonth=false}
\usepackage{lipsum}

\usepackage{marvosym}

\renewcommand*\sfdefault{lmss}
\renewcommand*\ttdefault{txtt}

% Information to be included in the title page:
\title{(Q)\texttt{M}-types and Coinduction in HoTT / CTT}
\subtitle{Master's Thesis, Computer Science}
\author{\large Lasse Letager Hansen, 201508114\\[1ex]{\small Supervisor: Bas Spitters}}
\institute{Aarhus University}
\date{\today}

\renewcommand{\insertshortauthor}{Lasse Letager Hansen}
\titlegraphic{
  \epsfig{file=logo.eps}
}

% Custom packages added:

% \usepackage{fullpage}
% \setlength\marginparwidth{0.7in}

\usepackage{todonotes} % \todo command

\usepackage{proof} % \infer command

\usepackage{tikz-cd}
\usepackage{xfrac} % \sfrac - sideways fraction


% \usepackage{mathtools} % multlined

% allow breaks in align and equation (a new page for example)
\allowdisplaybreaks

\renewcommand{\thefootnote}{\fnsymbol{footnote}} %footnote style

\usepackage{amsthm}

% \newtheoremstyle{plain}
% {\abovedisplayskip}   % ABOVESPACE
% {\belowdisplayskip}   % BELOWSPACE
% {\itshape}  % BODYFONT
% {0pt}       % INDENT (empty value is the same as 0pt)
% {\bfseries} % HEADFONT
% {.}         % HEADPUNCT
% {5pt plus 1pt minus 1pt} % HEADSPACE
% {}          % CUSTOM-HEAD-SPEC
\theoremstyle{plain} % default
\newtheorem{thm}{Theorem}[section]
\newtheorem{lem}[thm]{Lemma}
\newtheorem{prop}[thm]{Proposition}
\newtheorem{cor}{Corollary}

\theoremstyle{definition}
\newtheorem{defn}[thm]{Definition}% [section]
\newtheorem{conj}{Conjecture}[section]
\newtheorem{exmp}{Example} % [section]

\theoremstyle{remark}
\newtheorem*{rem}{Remark}
% \newtheorem*{note}{Note} % Not used?
\newtheorem{case}{Case}

% Custom commands

\newcommand*{\thmref}[1]{\hyperref[thm:#1]{Theorem~\ref*{thm:#1}}} % Theorem~\ref
\newcommand*{\lemref}[1]{\hyperref[lem:#1]{Lemma~\ref*{lem:#1}}}
\newcommand*{\corref}[1]{\hyperref[cor:#1]{Corollary~\ref*{cor:#1}}}
\newcommand*{\defref}[1]{\hyperref[defn:#1]{Definition~\ref*{defn:#1}}}
\newcommand*{\figref}[1]{\hyperref[fig:#1]{Figure~\ref*{fig:#1}}}
\newcommand*{\tableref}[1]{\hyperref[table:#1]{Table~\ref*{table:#1}}}
\newcommand*{\exmpref}[1]{\hyperref[exmp:#1]{Example~\ref*{exmp:#1}}}
\newcommand*{\sectionref}[1]{\hyperref[sec:#1]{Section~\ref*{sec:#1}}}
\newcommand*{\subsectionref}[1]{\hyperref[sec:#1]{Subsection~\ref*{sec:#1}}}
\newcommand*{\chapterref}[1]{\hyperref[ch:#1]{Chapter~\ref*{ch:#1}}}

\newcommand*{\set}[1]{\left\{#1\right\}}
\newcommand*{\term}[1]{\textcolor{green!30!black}{#1}} % \textcolor{green!70!black} \textcolor{blue!80!black}
\newcommand*{\pathterm}[1]{\textcolor{green!65!red!75!black}{#1}}
\newcommand*{\type}[1]{\textcolor{magenta!90!black}{#1}}
% \newcommand*{\dependenttype}[1]{\textcolor{magenta!70!black}{#1}}
\newcommand*{\container}[1]{\textcolor{orange}{#1}}
\newcommand*{\containerpair}[2]{\ensuremath{\colorlet{savedleftcolor}{.}\color{orange}\left(\color{savedleftcolor}#1\,\textcolor{orange}{\mathbf{,}}\,#2\color{orange}\right)\color{savedleftcolor}}}
\newcommand*{\containerpairsimple}[2]{\containerpair{#1}{\lambda\,\_,\,#2}}
\newcommand*{\universe}[1]{\textcolor{orange!80!black}{#1}}
\newcommand*{\unit}{\type{\textbf{1}}}
\newcommand*{\empt}{\type{\textbf{0}}}
\newcommand*{\coalg}[2]{#1\texttt{-}#2}

\newcommand*{\relation}[1]{\textcolor{gray!75!black}{\ensuremath{\mathtt{#1}}}}
\newcommand*{\constant}[1]{\textcolor{orange!60!black}{\ensuremath{\mathtt{#1}}}}

\newcommand*{\function}[1]{\textcolor{blue!60!black}{\ensuremath{\mathtt{#1}}}}
\newcommand*{\constructor}[1]{\textcolor{purple!60!black}{\ensuremath{\mathtt{#1}}}}
\newcommand*{\destructor}[1]{\textcolor{yellow!60!black}{\ensuremath{\mathtt{#1}}}}
\newcommand*{\typeformer}[1]{\ensuremath{\mathtt{#1}}}
\newcommand*{\functor}[1]{\ensuremath{\mathbf{\mathtt{#1}}}}

\newcommand*{\unitelem}{\constant{\star}} % \constant{tt}
\newcommand*{\quotientconstructor}[1]{\constructor{[}\,#1\,\constructor{]}}

\newcommand*{\ssfrac}[2]{ {\LARGE \text{\sfrac{\ensuremath{#1}}{\ensuremath{#2}}}}}

\newcommand*{\natcases}[2]{\ensuremath{\mathbb{[\hspace{-.5mm}\rangle}~\ensuremath{#1}~,~\ensuremath{#2}~\mathbb{\langle\hspace{-.5mm}]}}}

% Concatenation of paths taken from HoTT book
\newcommand{\ct}{%
  \mathchoice{\mathbin{\raisebox{0.5ex}{$\displaystyle\centerdot$}}}%
  {\mathbin{\raisebox{0.5ex}{$\centerdot$}}}%
  {\mathbin{\raisebox{0.25ex}{$\scriptstyle\,\centerdot\,$}}}%
  {\mathbin{\raisebox{0.1ex}{$\scriptscriptstyle\,\centerdot\,$}}}
}

\newcommand*{\sym}[1]{\ensuremath{#1^{-1}}}

\newcommand{\setlengths}{
  \setlength{\abovedisplayskip}{4pt}
  \setlength{\belowdisplayskip}{4pt}
  \setlength{\abovedisplayshortskip}{2pt}
  \setlength{\belowdisplayshortskip}{2pt}
}
\newcommand{\startwitheq}{\vspace{-2.5mm}} % -\baselineskip

\begin{document}

\frame{\titlepage}

\section{Introduction}

\begin{frame}
  \frametitle{Overview}
  \tableofcontents[subsubsectionstyle=hide]
\end{frame}

\subsection{Goals}
\begin{frame}
  \frametitle{Goals}
  \begin{itemize}
  \item Formalize coinductive types as \texttt{M}-types
  \item Define equality for coinductive types % given a bisimulation relation
  \item Explore ways to define quotiented \texttt{M}-types
  \end{itemize}
\end{frame}

\section{\texttt{M}-types}
\subsection{Definition of \texttt{M}-types}
% \begin{frame}[fragile]
%   \frametitle{\texttt{M}-types}
%   \begin{itemize}
%   \item \texttt{M}-types are non-wellfounded trees
%     \begin{figure}
%       \centering
%       \begin{tikzcd}[row sep=large]
%         & & \type{C} & \type{A} & \type{A} & \type{C} & \dots \\
%         & \type{C} & & \type{A} \ar[ul,"\typeformer{B}~\term{a}_0" description] \ar[u,"\typeformer{B}~\term{a}_1" description] \ar[ur,"\typeformer{B}~\term{a}_2" description] & & \type{A} \ar[u,"\typeformer{B}~\term{a}_0" description] \ar[ur,"\dots" description] \\
%         & & & \type{A} \ar[ull,"\typeformer{B}~\term{a}_0" description] \ar[u,"\typeformer{B}~\term{a}_1" description] \ar[urr,"\typeformer{B}~\term{a}_2" description]
%       \end{tikzcd}
%     \end{figure}
%   \item \texttt{M}-types are final coalgebras
%     \begin{figure}[h]
%       \centering
%       \begin{tikzcd}[row sep=large,column sep=large]
%         \type{C} \ar[d] \ar[r,dashed,"\function{f}"] & \typeformer{M} \ar[d,"\function{out}"] \\
%         \functor{P}\,\type{C} \ar[r,"\functor{P}\,\function{f}"] & \functor{P}\typeformer{M}
%       \end{tikzcd}
%       \caption{Final \functor{P}-Coalgebra}
%     \end{figure}
%   \end{itemize}
% \end{frame}

\begin{frame}[fragile]
  \frametitle{Containers and Polynomial functors}
  \begin{defn}\setlengths
    A Container (or signature) is a dependent pair \(\container{S} = \containerpair{\type{A}}{\typeformer{B}}\) for the types \(\type{A} : \universe{\mathcal{U}}\) and \(\typeformer{B} : \type{A} \rightarrow \universe{\mathcal{U}}\). 
  \end{defn}
  \begin{figure}
    \centering
    \begin{tikzcd}[row sep=large]
      \type{C} & \type{D} & \type{E} \\
      & \type{A} \ar[ul,"\typeformer{B}~\term{a}_0" description] \ar[u,"\typeformer{B}~\term{a}_1" description] \ar[ur,"\typeformer{B}~\term{a}_2" description]
    \end{tikzcd}
  \end{figure}
\end{frame}

\begin{frame}[fragile]
  \frametitle{Containers and Polynomial functors}
  \begin{defn}\setlengths
    A polynomial functor \(\functor{P}_{\container{S}}\) (or extension) for a container \(\container{S} = \containerpair{\type{A}}{\typeformer{B}}\) is defined, for types as
    \begin{equation}
      \functor{P}_{\container{S}}\,\type{X} = \sum_{(\term{a} : \type{A})} \typeformer{B}\term{a} \rightarrow \type{X}
    \end{equation}
    and for a function \(\function{f} : \type{X} \rightarrow \type{Y}\) as
    \begin{equation}
      \begin{aligned}
        \functor{P}_{\container{S}}\,\function{f}~(\term{a},\function{g}) = (\term{a}, \function{f} \circ \function{g}).
      \end{aligned}
    \end{equation}
  \end{defn}
\end{frame}

\begin{frame}[fragile]
  \frametitle{Chain}
  \begin{defn}[Chain]\setlengths
    We define a chain as a family of morphisms \(\function{\pi}_{(n)} : \type{X}_{n+1} \rightarrow \type{X}_{n}\), over a family of types \(\type{X}_{n}\). See figure.
  \end{defn}
  \begin{figure}[h]
    \centering
    \begin{tikzcd}[row sep=large,column sep=large]
      \type{X}_0 & \type{X}_1 \ar[l,"\function{\pi}_{(0)}"] & \cdots \ar[l,"\function{\pi}_{(1)}"] & \type{X}_n \ar[l,"\function{\pi}_{(n-1)}"] & \type{X}_{n + 1} \ar[l,"\function{\pi}_{(n)}"] & \cdots \ar[l,"\function{\pi}_{(n+1)}"]
    \end{tikzcd}
  \end{figure}
  \begin{defn}\setlengths
    The limit of a chain is given as
    \begin{equation}
      \type{\mathcal{L}} = \sum_{(\term{x} : \prod_{(n : \mathbb{N})} \type{X}_n)} \prod_{(n : \mathbb{N})} (\function{\pi}_{(n)}~\term{x}_{n+1} \equiv \term{x}_n)
    \end{equation}
  \end{defn}
  We let \(\typeformer{M}_{\container{S}}\) be the limit, for a chain defined by \(\type{X}_n = \functor{P}^n\,\unit\), and \(\function{\pi}_{(n)} = \functor{P}^n\,\constant{!}\)
\end{frame}

\begin{frame}
  \frametitle{Equality between \(\mathcal{L}\) and \(\functor{P}\,\mathcal{L}\)}
  \begin{thm}\setlengths
    There is an equality
    \begin{equation}
      \pathterm{shift} : \typeformer{M} \equiv \functor{P}\,\typeformer{M}
    \end{equation}
    from which we can define helper functions
    \begin{equation}
      \function{in} : \functor{P}\,\typeformer{M} \rightarrow \typeformer{M} \qquad \function{out} : \typeformer{M} \rightarrow \functor{P}\,\typeformer{M}
    \end{equation}    
  \end{thm}
  \pause
  \begin{block}{Proof structure}\setlengths
    The proof is done using the two helper lemmas
    \begin{align}
      \pathterm{\alpha} : \type{\mathcal{L}}^{\functor{P}} \equiv \functor{P}\,\type{\mathcal{L}} \\
      \pathterm{\mathcal{L}unique} : \type{\mathcal{L}} \equiv \type{\mathcal{L}}^{\functor{P}}
    \end{align}
    where \(\type{\mathcal{L}}^{\functor{P}}\) is the limit of the shifted chain defined as \(\type{X'}_n = \type{X}_{n+1}\) and \(\function{\pi'}_{(n)} = \function{\pi}_{(n+1)}\). With these two lemmas we get \(\pathterm{shift} = \pathterm{\alpha} \ct \pathterm{\mathcal{L}unique}\).
  \end{block}
\end{frame}

\begin{frame}
  \frametitle{Coalgebra}
  \framesubtitle{\texttt{M}-types are final coalgebras}
  We want to show that \texttt{M}-types are final coalgebras.
  \begin{defn}\setlengths
    A \(\functor{P}\)-coalgebra is defined as
    \begin{equation}
      \sum_{(\type{C} : \universe{\mathcal{U}})} \type{C} \rightarrow \functor{P}\,\type{C}
    \end{equation}
    which we denote \(\coalg{\type{C}}{\function{\gamma}}\). We define \(\functor{P}\)-coalgebra morphisms as
    \begin{equation}
      \coalg{\type{C}}{\function{\gamma}} \Rightarrow \coalg{\type{D}}{\function{\delta}} = \sum_{(\function{f} : \type{C} \rightarrow \type{D})} \function{\delta} \circ \function{f} \equiv \functor{P} \function{f} \circ \function{\gamma}
    \end{equation}
  \end{defn}
  A coalgebra is final, if the following is true
  \begin{equation}
    \sum_{(\coalg{\type{D}}{\function{\rho}})}\prod_{(\coalg{\type{C}}{\function{\gamma}})} \mathtt{isContr}~(\coalg{\type{C}}{\function{\gamma}} \Rightarrow \coalg{\type{D}}{\function{\rho}})    
  \end{equation}
\end{frame}

\begin{frame}
  \frametitle{\texttt{M}-types are final coalgebras}
  \begin{thm}\setlengths
    \texttt{M}-types are final coalgebras. That is \(\typeformer{Final}_{\container{S}}~\typeformer{M}\).
  \end{thm}
  \pause
  \begin{block}{Proof structure}\setlengths
    The definition of finality is
    \begin{equation}
      \prod_{(\coalg{\type{C}}{\function{\gamma}} : \mathtt{Coalg}_{\container{S}})} \mathtt{isContr}~(\coalg{\type{C}}{\function{\gamma}} \Rightarrow \coalg{\typeformer{M}}{\function{out}})
    \end{equation}
    which we show by \((\coalg{\type{C}}{\function{\gamma}} \Rightarrow \coalg{\typeformer{M}}{\function{out}}) \equiv \unit\).
  \end{block}
\end{frame}

\subsection{Construction of \texttt{M}-types and Examples}
\begin{frame}
  \frametitle{Overview}
  \tableofcontents[currentsubsection]
\end{frame}

\begin{frame}[fragile]
  \frametitle{Example: Delay Monad}
  A delay monad is defined by the two constructors\\
  \strut
  \hfill
  \begin{minipage}[h]{0.35\linewidth}
    \begin{equation}
      \vcenter{\infer{\constructor{now}~\term{r} : \typeformer{Delay}~\type{R}}{\term{r} : \type{R}}}
    \end{equation}
  \end{minipage}
  \hfill
  \begin{minipage}[h]{0.4\linewidth}
    \begin{equation}
      \vcenter{\infer{\constructor{later}~\term{t} : \typeformer{Delay}~\type{R}}{\term{t} : \typeformer{Delay}~\type{R}}}
    \end{equation}
  \end{minipage}
  \hfill
  \strut
  \\[2mm]
  We define a container
  \begin{equation}
    \containerpair{\type{R} + \unit}{[\empt, \unit]}
  \end{equation}
  and a polynomial functor
  \begin{equation}
    \functor{P}\,\type{X} = \quad \sum_{(\term{x} : \type{R} + \unit)} \begin{cases} \empt & \term{x} = \mathtt{inl}~\term{r} \\ \unit & \term{x} = \mathtt{inr}~\unitelem \end{cases} \rightarrow \type{X} \quad = \type{R} + \type{X},
  \end{equation}
  such that we get the diagram
  \begin{figure}[h]
    \centering
    \begin{tikzcd}[row sep=large,column sep=large]
      \type{R} \ar[dr,"\constructor{now}"'] \ar[r,"\texttt{inl}"] & \type{R} + \typeformer{Delay}~\type{R} \ar[d,shift left,"\function{in}"] & \typeformer{Delay}~\type{R} \ar[l,"\texttt{inr}"'] \ar[dl,"\constructor{later}"] \\
      & \typeformer{Delay}~\type{R} \ar[u,shift left,"\function{out}"]
    \end{tikzcd}
  \end{figure}
\end{frame}

\begin{frame}
  \frametitle{Rules for Constructing \texttt{M}-types}
  Adding containers \(\containerpair{\type{A}}{\typeformer{B}}\) and \(\containerpair{\type{C}}{\typeformer{D}}\) for two constructors together is done by
  \begin{equation}
    \containerpair{\type{A} + \type{C}}{ \begin{cases} \typeformer{B}\,\term{a} & \mathtt{inl}~\term{a} \\ \typeformer{D}\,\term{c} & \mathtt{inr}~\term{c} \end{cases}}
  \end{equation}
  whereas adding containers for two destructors is done by
  \begin{equation}
    \containerpair{\type{A} \times \type{C}}{\lambda\,(\term{a},\term{c}),~\typeformer{B}\,\term{a} + \typeformer{D}\,\term{c}}  
  \end{equation}
  However combining both destructors and constructors is not as simple. Which is similar to rules for coinductive records.
\end{frame}

\begin{frame}
  \frametitle{\texttt{M}-types and Records}
  We can take a coinductive record and transform it to an \texttt{M}-type. The types of fields in a coinductive records are
  \begin{itemize}
  \item non-dependent fields
  \item dependent fields
  \item recursive fields
  \item dependent and recursive fields
  \end{itemize}
  We will give some examples of these.
\end{frame}

\begin{frame}
  \frametitle{\texttt{M}-types and Records}
  \framesubtitle{Example: Record to M-type}
  Lets try and convert the record
  \begin{equation}
    \begin{aligned}
      &\mathtt{record}~\type{tree} : \universe{\mathcal{U} }~\mathtt{where} \\
      &\qquad \term{value} : \mathbb{N} \\
      &\qquad \term{left\text{-}child} : \type{tree} \\
      &\qquad \term{right\text{-}child} : \type{tree}
    \end{aligned}
  \end{equation}
  To an \texttt{M}-type defined by the container
  \begin{equation}
    \containerpair{\mathbb{N}}{\unit + \unit}
  \end{equation}
\end{frame}

\begin{frame}
  \frametitle{\texttt{M}-types and Records}
  \framesubtitle{Example: Record to M-type}
  Lets try and convert the record
  \begin{equation}
    \begin{aligned}
      &\mathtt{record}~\type{bet} : \universe{\mathcal{U} }~\mathtt{where} \\
      &\qquad \term{value_a} : \mathbb{N} \\
      &\qquad \term{value_b} : \mathbb{N} \\
      &\qquad \term{winner} : \term{value_a} \leq \term{value_b} \rightarrow \type{bool} \\
    \end{aligned}
  \end{equation}
  To an \texttt{M}-type defined by the container
  \begin{equation}
    \containerpair{\sum_{(\term{value_a} : \mathbb{N})} \sum_{(\term{value_b} : \mathbb{N})} \term{value_a} \leq \term{value_b} \rightarrow \type{bool}}{\empt}
  \end{equation}
\end{frame}


\begin{frame}
  \frametitle{\texttt{M}-types and Records}
  \framesubtitle{Example: Record to M-type}
  Lets try and convert the record
  \begin{equation}
    \begin{aligned}
      &\mathtt{record}~\typeformer{example}~\type{A} : \universe{\mathcal{U} }~\mathtt{where} \\
      &\qquad \term{value} : \type{A} \\
      &\qquad \type{index\text{-}type} : \universe{\mathcal{U}} \\
      &\qquad \term{continue} : \type{index\text{-}type} \rightarrow \typeformer{example}~\type{A} \\
    \end{aligned}
  \end{equation}
  To an \texttt{M}-type defined by the container
  \begin{equation}
    \containerpair{\type{A} \times \universe{\mathcal{U}}}{\lambda\,(\_,\type{index\text{-}type}), \type{index\text{-}type}}
  \end{equation}
\end{frame}

% \begin{frame}
%   \frametitle{\texttt{M}-types and Records}
%   We can take a coinductive record and transform it to an \texttt{M}-type by
%   \begin{itemize}
%   \item Adding all fields that are not a dependent \(\term{f}_1 : \type{F}_1,\term{f}_2 : \type{F}_2,\dots\) to the first part of the container as a product
%   \item Adding all fields that are a dependent \(\term{d}_1 : \type{D}_1, \term{d}_2 : \type{D}_2, \dots\) to the first part of the container as a dependent product
%   \item Adding all the fields that depend on other fields that depends on previously defined fields \(\term{b}_1 : \typeformer{B}_1~\term{d}_1, \dots\)
%   \item Converting all recursive fields to the last part of the container \(\term{r}_1 : \type{R}_1 \rightarrow \typeformer{M},\dots\)
%   \item Converting all dependent recursive fields to the last part of the container \(\term{z}_1 : \prod_{\term{v}_1 : \type{V}_1} \typeformer{Z}_1~\term{v}_1 \rightarrow \typeformer{M},\dots\)
%   \end{itemize}
%   \begin{equation}
%     \containerpair{\type{F}_1 \times .. \times \sum_{\term{d}_1 : \type{D}_1} \typeformer{B}_1~\term{d}_1 \times \sum_{\term{d}_2 : \type{D}_2} ..}{\lambda\,\_,\type{R}_1 \times .. \times \sum_{\term{v}_1 : \type{V_1}} \typeformer{Z}_1~\term{v_1} \times ..}
%   \end{equation}
% \end{frame}

% \begin{frame}[fragile]
%   \frametitle{\texttt{M}-types and Records}
%   \begin{figure}
%     \centering
%     \begin{tikzcd}[row sep=large,column sep=large]
%       &  & \type{F}_1 \ar[dl] \ar[d]  \\
%       \sum_{\term{d}_1 : \type{D}_1} \typeformer{B}_1 \term{d}_1 \ar[r] \ar[d] & \functor{P}\,\typeformer{M} & \dots \ar[d] \ar[l,draw=none,near start,"\vdots" description] \\
%       \type{D}_1 \ar[ur] & \typeformer{M}  & \type{F}_n \ar[ul]
%     \end{tikzcd}
%   \end{figure}
% \end{frame}

\subsection{Equality for Coinductive types}
\begin{frame}
  \frametitle{Overview}
  \tableofcontents[currentsubsection]
\end{frame}

\begin{frame}[fragile]
  \frametitle{Example: Streams}
  We can now define streams for a given type \(\type{A}\), as a records
  \begin{equation}
    \begin{aligned}
      &\mathtt{record}~\typeformer{Stream}~\type{A} : \universe{\mathcal{U} }~\mathtt{where} \\
      &\qquad \term{hd} : \type{A} \\
      &\qquad \term{tl} : \typeformer{Stream}~\type{A} \\
    \end{aligned}
  \end{equation}  
  corresponding to the container
  \begin{equation}
    \containerpair{\type{A}}{\unit}
  \end{equation}
  for which we get the polynomial functor
  \begin{equation}
    \functor{P}\,\type{X} = \type{A} \times \type{X}
  \end{equation}
  For which the we get the \texttt{M}-type for stream
  \begin{figure}[h]
    \centering
    \begin{tikzcd}
      \type{A} & \type{A} \times \typeformer{Stream}~\type{A} \ar[d,shift left,"\function{in}"] \ar[l,"\pi_1"'] \ar[r,"\pi_2"] & \typeformer{Stream}~\type{A} \\
      & \typeformer{Stream}~\type{A} \ar[u,shift left,"\function{out}"] \ar[ul,"\destructor{hd}"] \ar[ur,"\destructor{tl}"']
    \end{tikzcd}
  \end{figure}
\end{frame}

\begin{frame}[fragile]
  \frametitle{Bisimulation for Streams}
  We can define an equivalence relation
  \begin{equation}
    \vcenter{\infer{\term{s}~\relation{\sim_{stream}}~\term{t}}{\destructor{hd}~\term{s} \equiv \destructor{hd}~\term{t} & \destructor{tl}~\term{s}~\relation{\sim_{stream}}~\destructor{tl}~\term{t}}}
  \end{equation}
  We can formalize this as a bisimulation. A (strong) bisimulation for a \(\functor{P}\)-coalgebra \(\coalg{\type{C}}{\function{\gamma}}\) is given by
  \begin{itemize}
  \item a relation \(\relation{\mathcal{R}} : \type{C} \rightarrow \type{C} \rightarrow \universe{\mathcal{U}}\)
  \item a type \(\type{\overline{\mathcal{R}}} = \sum_{(\term{a} : \type{C})} \sum_{(\term{b} : \type{C})} \term{a}~\relation{\mathcal{R}}~\term{b}\)
  \item and a function \(\function{\alpha}_{\relation{\mathcal{R}}} : \type{\overline{\mathcal{R}}} \rightarrow \functor{P}_{\container{S}}\,\type{\overline{\mathcal{R}}}\)
  \end{itemize}
  Such that \(\coalg{\type{\overline{\mathcal{R}}}}{\function{\alpha}_{\relation{\mathcal{R}}}}\) is a \(\functor{P}\)-coalg, making the following diagram commute.
  \begin{figure}[h]
    \centering
    \begin{tikzcd}[row sep=huge,column sep=huge]
      \coalg{\type{C}}{\function{\gamma}} & \coalg{\type{\overline{\mathcal{R}}}}{\function{\alpha}_{\relation{\mathcal{R}}}} \ar[l,Rightarrow,"{\pi_1}^{\type{\overline{\mathcal{R}}}}"'] \ar[r,Rightarrow,"{\pi_2}^{\type{\overline{\mathcal{R}}}}"] & \coalg{\type{C}}{\function{\gamma}}
    \end{tikzcd}
  \end{figure}
  In MLTT this is not enough to define an equality, however in HoTT and CTT it is.
\end{frame}

\begin{frame}[fragile]
  \frametitle{Coinduction Principle}
  \begin{thm}[Coinduction principle]\setlengths
    Given a relation \(\relation{\mathcal{R}}\), that is a bisimulation for a \texttt{M}-type, then (strongly) bisimilar elements \(\term{x}~\relation{\mathcal{R}}~\term{y}\) are equal \(\term{x} \equiv \term{y}\).
  \end{thm}
  \pause
  \begin{proof}\setlengths
    We get the diagram
    \begin{figure}[h]
      \centering
      \begin{tikzcd}[row sep=huge,column sep=huge]
        \coalg{\typeformer{M}}{\function{out}} & \coalg{\type{\overline{\mathcal{R}}}}{\function{\alpha}_{\relation{\mathcal{R}}}} \ar[l,Rightarrow,"{\pi_1}^{\type{\overline{\mathcal{R}}}}"'] \ar[r,Rightarrow,"{\pi_2}^{\type{\overline{\mathcal{R}}}}"] & \coalg{\typeformer{M}}{\function{out}}
      \end{tikzcd}
    \end{figure}
    Since \(\coalg{\typeformer{M}}{\function{out}}\) is a final coalgebra, functions into it are unique, meaning
    \begin{equation}
    \pi_1^{\type{\overline{\mathcal{R}}}} \equiv \pi_2^{\type{\overline{\mathcal{R}}}} % \equiv ~\constant{!}
  \end{equation}
    therefore given \(\term{r} : \term{x}~\relation{\mathcal{R}}~\term{y}\), we can construct the equality
    \begin{equation}
      \term{x} \equiv \pi_1^{\type{\overline{\mathcal{R}}}} (\term{x}, \term{y}, \term{r}) \equiv \pi_2^{\type{\overline{\mathcal{R}}}} (\term{x}, \term{y}, \term{r}) \equiv \term{y}.
    \end{equation}
  \end{proof}  
\end{frame}

\section{Quotient \texttt{M}-types}
\begin{frame}
  \frametitle{Overview}
  \tableofcontents[currentsection]
\end{frame}

\subsection{Propositional Truncation and Set Truncated Quotient}
\begin{frame}
  \frametitle{Propositional Truncation and Set Truncated Quotient}
  A Higher Inductive Type (HIT) is a type defined by point constructors as well as equality constructors. \pause % We define propositional truncation by
  % \begin{defn}[Propositional Truncation]\setlengths\startwitheq
  %   \strut
  %   \hfill
  %   \begin{minipage}[b]{0.25\linewidth}
  %     \begin{equation}
  %       \vcenter{\infer{\constructor{|}\,\term{x}\,\constructor{|} : \| \type{A} \|}{\term{x} : \type{A}}}
  %     \end{equation}
  %   \end{minipage}
  %   \hfill
  %   \begin{minipage}[b]{0.4\linewidth}
  %     \begin{equation}
  %       \vcenter{\infer{\constructor{squash}~\term{x}~\term{y} : \term{x} \equiv \term{y}}{\term{x}, \term{y} : \|\type{A}\|}}
  %     \end{equation}
  %   \end{minipage}
  %   \hfill
  %   \strut
  % \end{defn}
  \pause
  We can define set truncated quotients as the following HIT.
  \begin{defn}[Set Truncated Quotient]\setlengths \startwitheq
    \strut
    \hfill
    \begin{minipage}[b]{0.25\linewidth}
      \begin{equation}
        \vcenter{\infer{\quotientconstructor{\term{x}} : \type{A} / \relation{\mathcal{R}}}{\term{x} : \type{A}}}
      \end{equation}
    \end{minipage}
    \hfill
    \begin{minipage}[b]{0.4\linewidth}
      \begin{equation}
        \vcenter{\infer{\constructor{eq/}~\term{x}~\term{y}~\term{r} : \quotientconstructor{\term{x}} \equiv \quotientconstructor{\term{y}}}{\term{x}, \term{y} : \type{A} & \term{r} : \term{x}~\relation{\mathcal{R}}~\term{y}}}
      \end{equation}
    \end{minipage}
    \hfill
    \strut
    \begin{equation}
      \vcenter{\infer{\constructor{squash/} :\texttt{isSet}~(\type{A} / \relation{R})}{\strut}}
    \end{equation}
  \end{defn}
\end{frame}

\begin{frame}
  \frametitle{Partiality monad}
  \framesubtitle{Q\texttt{M}-type}
  The partiality monad describes equality of partial computations.
  \pause
  \\[3mm]
  We define the partiality monad by quotienting the delay monad by a relation defined by the constructors
  \\
  \strut
  \hfill
  \begin{minipage}{0.4\linewidth}
    \begin{equation}
      \vcenter{\infer[{\constructor{\sim_{later_l}}}]{\constructor{later}~\term{x}~\relation{\sim}~\term{y}}{\term{x}~\relation{\sim}~\term{y}}}
    \end{equation}
  \end{minipage}
  \hfill
    \begin{minipage}{0.4\linewidth}
    \begin{equation}
      \vcenter{\infer[{\constructor{\sim_{later_r}}}]{\term{x}~\relation{\sim}~\constructor{later}~\term{y}}{\term{x}~\relation{\sim}~\term{y}}}
    \end{equation}
  \end{minipage}
  \hfill
  \strut
  \\
  \strut
  \hfill
  \begin{minipage}{0.4\linewidth}
    \begin{equation}
      \vcenter{\infer[{\constructor{\sim_{now}}}]{\constructor{now}~\term{a}~\relation{\sim}~\constructor{now}~\term{b}}{\term{a}~\equiv~\term{b}}}
    \end{equation}
  \end{minipage}
  \hfill
  \begin{minipage}{0.5\linewidth}
    \begin{equation}
      \vcenter{\infer[{\constructor{\sim_{later}}}]{\constructor{later}~\term{x}~\relation{\sim}~\constructor{later}~\term{y}}{\term{x}~\relation{\sim}~\term{y}}}
    \end{equation}
  \end{minipage}
  \hfill
  \strut
  \\[5mm]
  the partiality monad is then given by the set truncated quotient
  \begin{equation}
    \typeformer{Delay}~\type{R}/\relation{\sim}
  \end{equation}
  however we need the axiom of (countable) choice
\end{frame}

\begin{frame}
  \frametitle{Partiality Monad}
  \framesubtitle{QIIT}
  A QIIT is a type that is defined at the same time as a relation and then set truncated.
  \\[3mm]
  \pause
  We can define the partiality monad as a QIIT with type constructors\\[-4mm]
  \strut
  \hfill
  \begin{minipage}{0.25\linewidth}
    \begin{equation}
      \vcenter{\infer{\type{R}_\bot : \universe{\mathcal{U}}}{\strut}}
    \end{equation}
  \end{minipage}
  \hfill
  \begin{minipage}{0.25\linewidth}
    \begin{equation}
      \vcenter{\infer{\constructor{\bot} : \type{R}_\bot}{\strut}}
    \end{equation}
  \end{minipage}
  \hfill
  \begin{minipage}{0.25\linewidth}
    \begin{equation}
      \vcenter{\infer{\constructor{\eta}~\term{a} : \type{R}_\bot}{\term{a} : \type{R}}}
    \end{equation}
  \end{minipage}
  \hfill
  \strut
  \\[3mm]
  and an ordering relation \((\cdot~\relation{\sqsubseteq_{\bot}}~\cdot)\) indexed twice over \(\type{R}_\bot\) \\
  \strut
  \hfill
  \begin{minipage}{0.3\linewidth}
    \begin{equation}
      \vcenter{\infer[\constructor{\sqsubseteq_{refl}}]{\term{x}~\relation{\sqsubseteq_{\bot}}~\term{x}}{\term{x} : \type{R}_\bot}}
    \end{equation}
  \end{minipage}
  \hfill
  \begin{minipage}{0.45\linewidth}
    \begin{equation}
      \vcenter{\infer[\constructor{\sqsubseteq_{trans}}]{\term{x}~\relation{\sqsubseteq_{\bot}}~\term{z}}{\term{x}~\relation{\sqsubseteq_{\bot}}~\term{y} & \term{y}~\relation{\sqsubseteq_{\bot}}~\term{z}}}
    \end{equation}
  \end{minipage}
  \hfill
  \strut
  \\[4mm]
  \begin{equation}
    \vcenter{\infer{\constructor{\alpha_{\bot}}~\term{p}~\term{q} : \term{x} \equiv \term{y}}{\term{x},\term{y} : \type{R}_\bot & \term{p} : \term{x}~\relation{\sqsubseteq_{\bot}}~\term{y} & \term{q} : \term{y}~\relation{\sqsubseteq_{\bot}}~\term{x}}}
  \end{equation}
  where \(\constructor{\bot}\) is the smallest element in the order
  \begin{equation}
    \vcenter{\infer[\constructor{\sqsubseteq_{never}}]{\constructor{\bot}~\relation{\sqsubseteq_{\bot}}~\term{x}}{\term{x} : \type{R}_\bot}}
  \end{equation}
\end{frame}

\begin{frame}
  \frametitle{Partiality Monad}
  \framesubtitle{QIIT}
  with an upper bound
  \begin{equation}
    \vcenter{\infer{\constructor{\bigsqcup}(\function{s},\function{b}) : \type{R}_\bot}{\function{s} : \mathbb{N} \rightarrow \type{R}_\bot & \function{b} : \prod_{(n : \mathbb{N})} \function{s}_n~\relation{\sqsubseteq_{\bot}}~\function{s}_{n+1} }}
  \end{equation}
  which gives a bound for a sequence
  \begin{equation}
    \vcenter{\infer{\prod_{(n : \mathbb{N})} \function{s}_n~\relation{\sqsubseteq_{\bot}}~\constructor{\bigsqcup}(\function{s},\function{b})}{\function{s} : \mathbb{N} \rightarrow \type{R}_\bot & \function{b} : \prod_{(n : \mathbb{N})} \function{s}_n~\relation{\sqsubseteq_{\bot}}~\function{s}_{n+1} }}
  \end{equation}
  and which is a least upper bound
  \begin{equation}
    \vcenter{\infer{\constructor{\bigsqcup}(\function{s},\function{b})~\relation{\sqsubseteq_{\bot}}~\term{x}}{\prod_{(n:\mathbb{N})} \function{s}_n~\relation{\sqsubseteq_{\bot}}~\term{x}}}
  \end{equation}
  and finally set truncated by the constructor \(\constructor{(-)_\bot\text{-}isSet}\)
\end{frame}

\begin{frame}
  \frametitle{Partiality Monad}
  \framesubtitle{Equality}
  To show these two definitions are equal we
  \begin{itemize}
  \item Define an intermediate type of sequences
    \begin{equation}
      \typeformer{Seq}~\type{A} = \sum_{(\function{s} : \mathbb{N} \rightarrow \type{A} + \unit)}  \prod_{(n : \mathbb{N})} \function{s}_n~\relation{\sqsubseteq_{\type{R} + \unit}}~\function{s}_{n+1}
    \end{equation}
  \item Show that the Delay monad is equal to this intermediate type
  \item Define weak bisimilarty for sequences, and show it respects the equality to the Delay monad
  \item Show that \(\typeformer{Seq}_{\type{R}}/\relation{\sim}\) is equal to the partiality monad, by
    \begin{itemize}
    \item Defining a function from \(\typeformer{Seq}_{\type{R}}/\relation{\sim}\) to \(\type{R}_{\bot}\)
    \item Show this function is injective and surjective
    \end{itemize}
  \end{itemize}
  However the proof for surjectivity requires the axiom of countable choice!
\end{frame}

\begin{frame}
  \frametitle{Simple Partialty Monad QIIT}
  This construction is quite involved, however doing a trivial construction\\[3mm]
  \strut\hfill
  \begin{minipage}{0.3\linewidth}
    \begin{equation}
      \vcenter{\infer{\constructor{now}~\term{a} : \type{R}_\bot}{a : \type{R}}}
    \end{equation}
  \end{minipage}
  \hfill
  \begin{minipage}{0.4\linewidth}
    \begin{equation}
      \vcenter{\infer{\constructor{later}~\term{t} : \type{R}_\bot}{\term{t} : \type{R}_\bot}}
    \end{equation}
  \end{minipage}
  \hfill\strut
  \begin{equation}
    \vcenter{\infer[\constructor{later_{\equiv}}]{\constructor{later}~\term{x} \equiv \term{y}}{\term{x} \equiv \term{y}}}
  \end{equation}
  would not give us the intended functionality, so we need to think more about how we define the QIIT.
  \\[3mm]
  We have to think about how we define a Q\texttt{M}-type.
\end{frame}

\begin{frame}
  \frametitle{What do we want from a quotient \texttt{M}-type?}
  \begin{itemize}
  \item We would like to be able to construct a quotient from an \texttt{M}-type and a relation.
  \item We should be able to lift constructors to the quotient without the axiom of choice.
  \item The type should be equal to the type defined by the set truncated quotient if we assume the axiom of choice.
  \end{itemize}
\end{frame}

\begin{frame}
  \frametitle{Alternative: Quotient Polynomial Functor (QPF)}
  We can define quotiented \texttt{M}-types from a quotient polynomial functor.
  \begin{defn}[Quotient Polynomial Functor]\setlengths
    We define a quotient polynomial functor (QPF), for types as
    \begin{equation}
      \functor{F}\,\type{X} =  \sum_{(\term{a} : \type{A})} ((\typeformer{B}\,\term{a} \rightarrow \type{X}) / \relation{\sim}_{\term{a}})
    \end{equation}
    and for a function \(\function{f} : \type{X} \rightarrow \type{Y}\), we use the quotient eliminator with
    \begin{equation}
      \typeformer{P} = \lambda\,\_, (\typeformer{B}\,\term{a} \rightarrow \type{Y})/\relation{\sim}_{\term{a}}
    \end{equation}
    for which we need \(\function{\sim_{ap}}\), which says that given a function \(\function{f}\) and \(\function{x}~\relation{\sim}_{\term{a}}~\function{y}\) then \(\function{f} \circ \function{x}~\relation{\sim}_{\term{a}}~\function{f} \circ \function{y}\).
    \begin{equation}
      \functor{F}\,\function{f}~(\term{a} , \function{g}) = (\term{a} , \function{elim}~\function{g})
    \end{equation}
  \end{defn}
\end{frame}

\begin{frame}[fragile]
  \frametitle{Alternative: Quotient Polynomial Functor (QPF)}
  We get the diagram
  \begin{figure}
    \centering
    \begin{tikzcd}
      \unit \ar[d] & \functor{P}\,\unit \ar[l,"\function{\pi}_{(1)}"'] \ar[d] & \dots \ar[l,"\function{\pi}_{(2)}"'] & \typeformer{M} \ar[d] \ar[l] & \functor{P}\,\typeformer{M}  \ar[d] \ar[l] \\
      \unit & \functor{F}\,\unit \ar[l,"\function{\pi'}_{(1)}"] & \dots \ar[l,"\function{\pi'}_{(2)}"] & \typeformer{QM} \ar[l] & \functor{F}\,\typeformer{QM} \ar[l]
    \end{tikzcd}
  \end{figure}
  For which \(\typeformer{M} \equiv \functor{P}\,\typeformer{M}\) and we would hope \(\typeformer{QM} \equiv \functor{F}\,\typeformer{QM}\), however this requires the axiom of choice.
  \\[3mm]
  \pause
  However looking further into this is future research.
\end{frame}

\section{Conclusion}
\begin{frame}
  \frametitle{Conclusion}
  We have
  \begin{itemize}
  \item given a formalization/semantic of \texttt{M}-types % TODO: formalized
  \item shown examples of and rules for how to construct \texttt{M}-types
  \item given a coinduction principle for \texttt{M}-types
  \item described the construction of the partiality monad as a QIIT
  \item discussed ways of constructing quotient \texttt{M}-types
  \end{itemize}
  Contribution
  \begin{itemize}
  \item Formalization in Cubical Agda
  \item Introducing Q\texttt{M}-types
  \end{itemize}
\end{frame}

\begin{frame}
  \frametitle{Future Work}
  \begin{itemize}
  \item Indexed \texttt{M}-types
  \item Showing finality of Q\texttt{M}-types and fully formalizing the constructions % TODO: (not just rules)
  \item Equality between Coinductive records and \texttt{M}-types
  \item Explore Guarded Cubical Type Theory
  \end{itemize}
\end{frame}

\end{document}

% TODO: What cannot be done with M-type, but need Indexed M-type (example)
% TODO: Where would GCTT be needed / might be better ?